% Copyright (c) 2018 by Daniel R. Driver.
% !TEX encoding = UTF-8 Unicode
% !TEX TS-program = XeLaTeX

\documentclass[titlepage]{article}

% This document presumes a file structure and set of inputs that are
% available at: git@github.com:danieldriver/syllabi.git

\newcommand\policy{../policy}
\newcommand\incl{../includes}
\ProvidesFile{variables.tex}[2022/09/11 v2.3 -- Syllabus variables]

\usepackage{xspace} % make manual spaces (like \mycmd\ ) unnecessary
\usepackage{xifthen} % provides \isempty test

% variables for internal use
\newcommand\prof{}
\newcommand\pdegree{}
\newcommand\pphone{}
\newcommand\pemail{}
\newcommand\poffice{}
\newcommand\phours{}
%
\newcommand\ccode{}
\newcommand\ctitle{}
\newcommand\cseries{}
\newcommand\cversion{}
\newcommand\csemester{}
\newcommand\cmeetson{}
\newcommand\cmeetsat{}
\newcommand\cmeetsin{}
\newcommand\cwebsite{}
\newcommand\cdescrip{}
\newcommand\cprereqs{}
\newcommand\edobject{}

% in case of fully online, BTh, or undergraduate courses - https://tex.stackexchange.com/a/5896
\newif\ifonline
\newcommand\Int[2]{\ifonline#1\else#2\fi}
\newif\ifbth
\newcommand\bth[2]{\ifbth#1\else#2\fi}
\newif\ifUG
\newcommand\UG[2]{\ifUG#1\else#2\fi}

% commands for setting variables in the preamble
\newcommand\professor[2][PhD]{
  \renewcommand\pdegree{#1\xspace}
  \renewcommand\prof{#2\xspace}}
\newcommand\phone[1]{
  \renewcommand\pphone{\addfontfeatures{Numbers=Monospaced}#1\xspace}}
\newcommand\email[1]{
  \renewcommand\pemail{\href{mailto:#1}{#1}\xspace}}
\newcommand\officehours[2][AST Library/Teams]{
  \renewcommand\poffice{#1\xspace}
  \renewcommand\phours{\href{https://outlook.office.com/bookwithme/user/b97b147a2e5c43b4a0e200c2cb630b77@astheology.ns.ca?anonymous}{#2}\xspace}%
  }
%
\newcommand\coursecode[2][1.0]{
  \renewcommand\cversion{#1\xspace}
  \renewcommand\ccode{#2\xspace}}
\newcommand\coursetitle[2][]{
  \ifthenelse{\isempty{#1}}%
    {}% do nothing if #1 is empty, else:
    {\renewcommand\cseries{#1\\[1ex]}}
  \renewcommand\ctitle{#2\xspace}}
\newcommand\semester[1]{
  \renewcommand\csemester{#1\xspace}}
\newcommand\meets[3]{
  \newcommand\AM{\textsc{am}}
  \newcommand\PM{\textsc{pm}}
  \renewcommand\cmeetson{#1\xspace}
  \renewcommand\cmeetsat{\Int{From 9:00 \AM}{#2}\xspace}
  \renewcommand\cmeetsin{\Int{\href{https://smu.brightspace.com/d2l/login}{Brightspace}}{#3}\xspace}}
\newcommand\website[2][]{
  \renewcommand\cwebsite{
  \ifthenelse{\isempty{#1}}%
    {\href{https://#2}}% copy url if #1 is empty, else
    {\href{#1}}% use specified url and print #2
  {#2}\xspace}}
\newcommand\cdescription[2][BF 1001; RM 1000 or GTRS 6000]{
  \renewcommand\cprereqs{#1}
  \renewcommand\cdescrip{#2\par}}
\newcommand\objectives[1]{
  \renewcommand\edobject{#1\par}}


%\onlinetrue % \Int{true}{false}
\coursecode[2.0]{HB 3114}
\coursetitle{The Twelve Prophets}

% Taught as RLGS 3143 in:
%   - Winter 2010
%   - Winter 2013
% Taught as HB 3114 in:
%   - Winter 2018

\professor{Daniel R. Driver}
\phone{902-425-7051}
\email{ddriver@astheology.ns.ca}
\officehours{Tuesdays, 2:00--4:00 \PM}

\semester{Winter Term 2018}
\meets{Mondays}% \meets{on}{at}{in}
      {9:00--11:30 \AM}
      {Classroom 3}
\website{astheology.ns.ca}
\cdescription{% copy from the current Academic Calendar
	Ben Sira writes of “the bones of the Twelve Prophets,” asking that
	they may “send forth new life from where they lie, for they
	comforted the people of Jacob and delivered them with confident
	hope” (Sirach 49:10). It is a tantalizingly early reference to the
	association of Twelve Prophets brought together in the Hebrew
	Bible’s last book of Prophets, a collection later known as The Book
	of the Twelve. It is also a pregnant anticipation of the manifold
	ways in which the words of these prophets have spoken far beyond
	their original circumstances.

	To better understand the origins and impact of The Book of the
	Twelve, this course will look at its redactional history and
	editorial shaping, its final canonical forms, and its effects as
	biblical prophecy. It will introduce current issues in scholarly
	debate about The Twelve while also exploring ways that The Twelve
	meet the ancient expectation that they can “send forth new life”
	amidst the people of God. In broad terms, the course offers an
	encounter with several individual texts embedded in a wider biblical
	tradition, and an opportunity to see how together “the goodly
	fellowship of the Prophets” witnesses to long-term works of
	providence as judgment and mercy converge in “the great and terrible
	Day of the Lord” (Malachi 4:5).
	\emph{Prerequisite: \textsc{BF} 1001.}
}% end of course description
%\objectives{% recall Bloom's taxonomy: http://www.celt.iastate.edu/teaching/RevisedBlooms1.html
%By the end of the course students should be able to:
%}% end of learning objectives

\ProvidesFile{preamble.tex}[2014/09/02 v1.1 -- Syllabus preamble]

% basic typography
\usepackage{fontspec,realscripts}
\defaultfontfeatures{Ligatures=TeX}
\setmainfont[%
  Numbers={Proportional,OldStyle},
  SmallCapsFeatures={RawFeature=+c2sc;+smcp,LetterSpace=1}]{Scala OT}
\newfontfamily\lining[Numbers={Lining}]{Scala OT}
\setsansfont[%
  SmallCapsFeatures={RawFeature=+c2sc;+smcp,LetterSpace=1}]{Scala Sans Pro}
\setmonofont[Scale=MatchLowercase]{Menlo}
\usepackage{sectsty}
\allsectionsfont{\sffamily\scshape}
\frenchspacing
\setlength{\emergencystretch}{3em} % prevent overfull lines

% custom font size and leading
\renewcommand\tiny{\fontsize{6}{9}\selectfont}
\renewcommand\scriptsize{\fontsize{7}{10}\selectfont}
\renewcommand\footnotesize{\fontsize{8}{10}\selectfont}
\renewcommand\small{\fontsize{9}{11}\selectfont}
\renewcommand\normalsize{\fontsize{10}{12}\selectfont}% base size
\renewcommand\large{\fontsize{11}{13}\selectfont}
\renewcommand\Large{\fontsize{12}{14}\selectfont}
\renewcommand\LARGE{\fontsize{16}{18}\selectfont}% "course syllabus \\ semester" benefits from more lead
\renewcommand\huge{\fontsize{20}{22}\selectfont}
\renewcommand\Huge{\fontsize{24}{26}\selectfont}

% layout packages: page, logo, tables
\usepackage[scale={0.6,0.8},
            xetex]{geometry}
\usepackage{graphicx}
\usepackage{array}     % allow insertions of column styling with >{}
\usepackage{booktabs}  % elegant horizontal rules in tables

% custom macros for a session count in the schedule of readings
\newcounter{session}
\newcounter{columns}
\newcounter{spancolumns}
\newcounter{midcolumns}
\newcounter{courseunit}
\newcommand\setcolumncount[2][0]{ % optionally set count to other than 0,
  \setcounter{session}{#1}        % e.g. to -1, or to a standing count
  \setcounter{columns}{#2}
  \setcounter{spancolumns}{#2-1}
  \setcounter{midcolumns}{#2-2}}
\newcommand\sessioncount{\stepcounter{session}\arabic{session}}
\newcommand\sessionskip[1]{\multicolumn{1}{@{}r@{ }}{#1}}
\newcommand\unit[1]{\multicolumn{\thecolumns}{c}{%
  \scshape\stepcounter{courseunit}\roman{courseunit}. #1}}
\newcommand\noclass[1]{\sessionskip{---.}&\multicolumn{\themidcolumns}{@{}l}{\itshape No Class: #1}}
\newcommand\rarr{\Large\char"EA31\hspace*{0.5pt}} % pointer finger w/ reminder tie in Scala OT
\newcommand\reminder[2]{\sessionskip{\rarr}&\multicolumn{\thespancolumns}{@{}l}{#1 \hfill #2}}

% color to match Tyndale's branding
\usepackage[usenames]{xcolor}
% predefined: black, white, red, green, blue, cyan, magenta, yellow
\definecolor{TyndaleURLs}{HTML}{0062A0} % links on tyndale.ca
\definecolor{TyndaleBlue}{cmyk}{1,1,0,.32}
\definecolor{TyndaleGold}{cmyk}{0,.27,1,0}
\definecolor{TyndaleRed}{cmyk}{0,1,.99,.04}
\definecolor{TyndaleBlack}{cmyk}{0,0,0,1}
\definecolor{TyndaleGreen}{cmyk}{.45,0,1,.24}
\definecolor{TyndaleOrange}{cmyk}{0,.79,1,0}
\definecolor{TyndaleAqua}{cmyk}{.47,0,.24,0}
\definecolor{TyndaleYellow}{cmyk}{.03,.03,.35,0}

% metadata (assumes a host of definitions are made in the main file)
\usepackage[setpagesize=false,     % leave this to geometry
            hyperfootnotes=false,  % fragile and distracting
            xetex]{hyperref}
\hypersetup{breaklinks=true,       % allow link text to break across lines
            colorlinks=true,       % colorlinks resets pdfborder to 0 0 0
            urlcolor=TyndaleURLs,  % for external links
            linkcolor=TyndaleRed,  % for normal internal links
            citecolor=TyndaleGold, % for bibliographical citations in text
            pdfauthor={\prof},
            pdftitle={\ccode: \ctitle},
            pdfsubject={Tyndale UC, \csemester},
            pdfcreator={github.com/danieldriver/syllabus}}
\urlstyle{same}                    % don't use monospace font for urls

% custom footlines
\usepackage{fancyhdr}
\pagestyle{fancy} % turn it on
\fancyhf{}        % reset everything
\renewcommand{\headrulewidth}{0pt} % remove header line as well
\lfoot{\sffamily\scshape\footnotesize\ctitle, \csemester}
\rfoot{\sffamily\scshape\footnotesize\prof\quad\thepage}

% language support: polyglossia needs to go last
\usepackage{polyglossia}
\setdefaultlanguage{english}
\setotherlanguage{hebrew}
\newfontfamily\hebrewfont[Script=Hebrew]{SBL BibLit}

% protect positioning of margin table in policy/grades -- needs to come after polyglossia
\usepackage{marginfix}

% gratuitous with custom title page, but useful as a fallback
\title{\ccode: \ctitle}
\author{\professor}
\date{\semester}


\begin{document}
\ProvidesFile{title.tex}[2013/09/06 v1.0 -- Syllabus title page]

\begin{titlepage}
  \begin{center}

    \begin{minipage}{\textwidth}
      \parbox[t]{0.5\textwidth}{
        \mbox{}\\[-\baselineskip] % dummy line to align parboxes
        \includegraphics[width=0.5\textwidth]{.syllabus/includes/TyndaleUC}}
      \hfill
      \parbox[t]{0.4\textwidth}{
        \raggedleft\LARGE\sffamily Course Syllabus\\
        \semester}
    \end{minipage}

    \vfill

    {\sffamily\textsc{\LARGE\MakeLowercase\coursecode}\\[2ex]
      \bfseries\itshape\Huge\coursetitle}

    \vfill

    \begin{tabular}{>{\bfseries}rl>{\bfseries}rl}
      \toprule
      Instructor & \professor                   & Course  & \liningnums{\coursecode}~\texttt{(v\version)} \\
      \midrule
      Phone      & \texttt{\phone}              & Meets   & \meeton                                       \\
      Email      & \href{mailto:\email}{\email} & Time    & \meetat                                       \\
      Office     & \officenumber                & Room    & \meetin                                       \\
      Hours      & \officehours                 & Website & \href{http://\website}{\website}              \\
      \bottomrule
    \end{tabular}

    \vfill

    \begin{description}
      \item[Commuter Hotline]
        Class cancellations due to inclement weather or illness will
        be announced on the commuter hotline at \texttt{416.226.6620
        x2187}. Alternately, weather cancellation information is posted
        at \href{http://tyndale.ca/weather}{tyndale.ca/weather}.
      \item[MyTyndale.ca]
        This course may have materials stored on its website, such as
        handouts or readings that may be needed in order to complete
        assignments. Students are responsible for checking these course
        pages on a regular basis. Here, too, students are able to view
        their grades throughout the semester. For more information see
        Section~\ref{mytyndale}, below.
      \item[Mail]
        Students are responsible for information communicated through
        their campus mailboxes and student e-mail accounts. A mailbox
        directory hangs beside the mailboxes. For more information
        contact the Registrar's office.
    \end{description}

  \end{center}

  \section{Course Description}
  \label{description}

  \emph{From the Academic Calendar:} \cdescription

\end{titlepage}
\setcounter{page}{2} % count the title page as page 1


%  \section{Learning Objectives}
%  \label{objectives}
%
%  \edobject

\section{Required Texts \& Materials}
\label{texts}

The following texts are required. Students are strongly encouraged to
purchase their own copies. \href{http://danieldriver.com/courses/hb-3114/#required-texts-winter-2018}{Links to help you order the correct editions} are on the professor's website.

\begingroup
\renewcommand{\section}[2]{}% temporarily remove the section heading
\begin{thebibliography}{Theodoret}% use the longest item in the bibliography

	\bibitem[Two Sides]{Two} Ehud Ben Zvi and James D. Nogalski.
	\emph{Two Sides of a Coin: Juxtaposing Views on Interpreting the Book of the Twelve / the Twelve Prophetic Books}.
	Piscataway, NJ: Gorgias Press, 2009.
%	\textsc{isbn} 978-1607243038.

	\bibitem[Seitz]{Seitz} Christopher R. Seitz.
	\emph{The Goodly Fellowship of the Prophets: The Achievement of Association in Canon Formation}.
	Grand Rapids: Baker Academic, 2009.
%	\textsc{isbn} 978-0801038839.

	\bibitem[Jerome]{Jerome1} Jerome; Thomas P. Scheck, trans.
	\emph{Commentaries on the Twelve Prophets, Vol. 1 (Ancient Christian Texts)}.
	Downers Grove: IVP Academic, 2016.
%	\textsc{isbn} 978-0830829163.

	\bibitem[Theodoret]{Theodoret} Theodoret; Robert C. Hill, trans.
	\emph{Theodoret of Cyrus: Commentary on the Prophets, Vol. 3: Commentary on the Twelve Prophets}.
	Brookline, MA: Holy Cross Orthodox Press, 2007.
%	\textsc{isbn} 978-1885652768.

\end{thebibliography}
\endgroup

If you do not have access to a good study Bible, I recommend either the
NRSV (Michael Coogan, ed.) or the NJPS (Adele Berlin and Marc Zvi
Brettler, eds.), both published by Oxford University Press.

\section{Supplementary Texts}
\label{supplementary}

The following titles may be referenced as supplementary texts. They are
not required. Other supplementary materials may be provided through the
course website.

\begingroup
\renewcommand{\section}[2]{}% temporarily remove the section heading
\begin{thebibliography}{3}% use the longest item in the bibliography

	\bibitem{Bickerman} Elias Bickerman.
	\emph{Four Strange Books of the Bible: Jonah, Daniel, Koheleth, Esther}.
	New York: Schocken Books, 1967.

	\bibitem{Hill} Hill, Robert C.
	\emph{Reading the Old Testament in Antioch}. Bible in Ancient Christianity 5.
	Leiden: Brill, 2005. Link: \href{http://ezproxy.astheology.ns.ca:2048/login?url=http://search.ebscohost.com/login.aspx?direct=true&db=nlebk&AN=173885&site=ehost-live&scope=site}{AST eBook}.

	\bibitem{Jerome2} Jerome; Thomas P. Scheck, ed.
	\emph{Commentaries on the Twelve Prophets, Vol. 2 (Ancient Christian Texts)}.
	Downers Grove: IVP Academic, 2017.

	\bibitem{Nogalski} James D. Nogalski.
	\emph{The Book of the Twelve (2 Volumes)}. Smyth \& Helwys Bible commentary.
	Macon, GA: Smyth \& Helwys, 2011.

	\bibitem{Nog17} James D. Nogalski.
	\emph{The Book of the Twelve and Beyond: Collected Essays of James D. Nogalski}.
	Atlanta: SBL, 2017.

	\bibitem{Nogalski-Sweeney} James D. Nogalski and Marvin A. Sweeney, eds.
	\emph{Reading and Hearing the Book of the Twelve}. SBL Symposium Series.
	Atlanta: SBL, 2000.

	\bibitem{Redditt-Schart} Paul L. Redditt and Aaron Schart, eds.
	\emph{Thematic Threads in the Book of the Twelve}.
	Berlin: Walter de Gruyter, 2003.

	\bibitem{Sweeney} Marvin A. Sweeney.
	\emph{The Twelve Prophets (2 Volumes)}. Berit Olam commentary.
	Collegeville, MN: Liturgical, 2000.

	\bibitem{Watts-House} James W. Watts and Paul R. House, eds.
	\emph{Forming Prophetic Literature: Essays on Isaiah and the Twelve in Honor of John D.W. Watts}.
	Sheffield: Sheffield Academic, 1996.

\end{thebibliography}
\endgroup

\section{Course Outline}
\label{outline}

We will adhere to the schedule in \autoref{schedule} as closely as
possible, though the professor reserves the right to adjust it to suit
the needs of the class.

\setcolumncount{3}% set up \sessioncount, \unit{}, \noclass{}, and \reminder{memo}{date} macros
\begin{table}[htbp]% add p to put the schedule on its own page
  \centering
  \begin{tabular}{>{\sessioncount.}r@{ }lr}% make sure the column config agrees with \setcolumncount
	\toprule
	\sessionskip{\textbf{\S.}}&\textbf{Seminar Text} (Commentaries follow suit)&\textbf{Date}\\
	\midrule
		& None: Introductions, overview of The XII &  9 Jan. \\
		& The entire Book of the XII, from Hosea to Malachi & 16 Jan. \\
		& Jonah 1–2                                         & 23 Jan. \\
		& Jonah 3–4                                         & 30 Jan. \\
	\reminder{Review Essay 1: Due by the end of Week Four}{2 Feb.} \\
		& Micah 1–4                                         &  6 Feb. \\
		& Micah 5–7                                         & 13 Feb. \\
	\noclass{Reading Week}                                  & 20 Feb. \\
		& Nahum 1–2                                         & 27 Feb. \\
		& None: Attend Grad Projects in lieu of regular class & 6 Mar. \\
	\reminder{Review Essay 2: Due by the end of Week Eight}{9 Mar.} \\
		& Nahum 3                                           & 13 Mar. \\
		& Habakkuk 1–2                                      & 20 Mar. \\
		& Habakkuk 3                                        & 27 Mar. \\
	\reminder{Final Paper: Due by the end of Week Eleven}{30 Mar.} \\
		& The entire Book of the XII, from Hosea to Malachi & 3 Apr. \\
	\bottomrule
  \end{tabular}
  \caption{Schedule of Readings}
  \label{schedule}
\end{table}

See the AST website for a list of other \href{http://www.astheology.ns.ca/students/academic-dates.html}{important dates}.

\section{Evaluation}
\label{evaluation}

\subsection{Grade Structure for \ccode}
\label{structure}

\begin{enumerate}

	\item Weekly \textbf{seminar discussions} will facilitate close work
	with text and interpretation in select books from The XII.  Students
	should come to class fully prepared to discuss the assigned
	material, which will ordinarily include one to three chapters of
	primary material per week, plus commentary on those chapters.
	Commentary includes \cite{Jerome1} and \cite{Theodoret} at a
	minimum, but will ideally be extended to other ancient and modern
	readers. The quality of participation will be assessed across the
	semester.

	\item A \textbf{first review essay}, on \cite{Two}, will be due at
	the end of the fourth week of class (Friday, 2 February 2018). It
	should be 2,500 words long, plus or minus 10\%. The first third of
	the essay should be a summary of the work, while the remaining two
	thirds should be devoted to critical engagement and evaluation.

	\item A \textbf{second review essay}, on \cite{Seitz}, will be due
	at the end of the eighth week of class (Friday, 9 March 2018). It
	should be 2,500 words long, plus or minus 10\%. The first third of
	the essay should be a summary of the work, while the remaining two
	thirds should be devoted to critical engagement and evaluation.

	\item A comprehensive \textbf{final paper} will give students an
	opportunity to explore one aspect of the Book of the XII, or a short
	passage from a prophetic book within The XII, in detail. The paper
	must have a thesis that involves original exposition of and
	theological reflection on a biblical text. It is due at the end of
	the eleventh week of class (Friday, 30 March 2018). It should be
	3,500 words long, plus or minus 10\% (the total count includes
	footnotes but not the final bibliography).

\end{enumerate}

The breakdown for the semester's total work is shown in
\autoref{grade-dist}.

\begin{table}[htbp]
  \centering
  {\lining
  \begin{tabular}{lr}
    \toprule
    Weekly Seminars & 25\% \\
    Review Essay 1  & 20\% \\
    Review Essay 2  & 20\% \\
    Final Paper     & 35\% \\
    \bottomrule
  \end{tabular}}
  \caption{Distribution of Grades}
  \label{grade-dist}
\end{table}

\ProvidesFile{grades.tex}[2023/09/06 v2.5 -- Course policy]

AST's \href{https://www.astheology.ns.ca/students/resources.html}{Academic
Calendar} provides guidelines and criteria for academic assessment. Marks are
assigned by letter grade using these benchmarks. Note that graduate and
undergraduate programs are currently assessed on a different scales at AST (see
\autoref{grade-syst}). For graduate students, a grade of C is deemed
``Unsatisfactory,'' whereas for undergraduate (BTh) students, a C is
``Satisfactory'' and the grade of D is possible for ``Marginal'' performance.
Other grades have different definitions, too.

\begin{table}[htbp]
 \centering
 {\lining\footnotesize
 \begin{tabular}{lll|ll}
   \toprule
               & \multicolumn{2}{c|}{Graduate Courses}   & \multicolumn{2}{c}{Undergraduate Courses} \\
               \cline{2-5}
   A+          & 94--100    & Exceptional               & 90--100    & Excellent                 \\
   A           & 87--93     & Outstanding               & 85--89     & Excellent                 \\
   A\char"2212 & 80--86     & Excellent                 & 80--84     & Excellent                 \\ [1ex]
   B+          & 77--79     & Good                      & 75--79     & Good                      \\
   B           & 73--76     & Acceptable                & 70--74     & Good                      \\
   B\char"2212 & 70--72     & Marginally Acceptable     &            &                           \\ [1ex]
   C           & 60--69     & Unsatisfactory            & 60--69     & Satisfactory              \\
   D           &            &                           & 50--59     & Marginal                  \\
   F           & 0--59      & Failure                   & 0--49      & Failure                   \\
   FP          & 0          & Failure due to Plagiarism & 0          & Failure due to Plagiarism \\
   \bottomrule
 \end{tabular}}
 \caption{Outline of Grades at AST}
 \label{grade-syst}
\end{table}

\ProvidesFile{other.tex}[2018/0/03 v2.4 -- Course policy]

\section{Other Course Policy}
\label{policy}

Late work will not be accepted, except in genuinely extenuating
circumstances. Students must submit something before the deadline if
they wish to receive credit. Unless I state otherwise, written
assignments are to be uploaded by 11:59 \PM\ on the date indicated.

Essay submissions must be typewritten, double-spaced, and formatted as
PDFs. They should be free from error. In this course they should follow
either \emph{The SBL Handbook of Style} (refer to the online
\href{https://www.sbl-site.org/assets/pdfs/pubs/SBLHSsupp2015-02.pdf}{
Student Supplement}), or \emph{The Chicago Manual of Style}, per AST's
policy for major papers. As a reminder, AST also upholds an Inclusive
Language Policy.

\href{http://www.eerdmans.com/Pages/Item/59043/Commentary-Statement.aspx}{Plagiarism},
if \href{https://www.theguardian.com/world/2013/feb/09/german-education-minister-quits-phd-plagiarism}{detected},
will result in failure of the course.

Students should request permission to record a class or lecture. If
permission is granted, or if recordings are provided (as in the case of
an online course), I stipulate that all recordings be for personal use
only. They may not be shared or distributed.

If you have abilities or disabilities that require modifications to the
assessment process or other aspects of this course, please advise the
course instructor as soon as possible.

Finally, I encourage the conscientious use of laptops, tablets, and
other technology in my classes. In classroom settings, realize that, as
\href{http://dx.doi.org/10.1016/j.compedu.2012.10.003}{cognitive
psychologists have demonstrated}, ``laptop multitasking hinders
classroom learning for both users and nearby peers.'' Do your part to
foster an environment of open dialogue by honouring the presence of your
classmates. In online settings, consider both the physical environment
in which you choose to work and the virtual environment that you help
create through your participation in various forums. Let your engagement
in this course be marked by rigour and charity alike.

\end{document}
